%%%%%%%%%%%%%%%%%%%%%%%%%%%%%%%%%%%%%%%%%%%%%%%%%%%%%%%%%%%%
% 12.4 PROOF THEORY
% The Composition of Advanced Mathematics
%%%%%%%%%%%%%%%%%%%%%%%%%%%%%%%%%%%%%%%%%%%%%%%%%%%%%%%%%%%%

\section{Proof Theory}
\label{sec:proof-theory}

\prerequisite{
Formal logic, propositional logic, first-order logic, natural deduction,
formal derivations, syntax and semantics, soundness, completeness,
mathematical induction, and elementary model theory.
}

\learningobjectives{
By the end of this section, the reader should be able to:
\begin{itemize}
    \item explain the purpose of proof theory;
    \item treat formal proofs as mathematical objects;
    \item distinguish axioms from inference rules;
    \item understand derivations and proof trees;
    \item distinguish object-level proofs from metatheoretic proofs;
    \item recognize Hilbert-style proof systems;
    \item understand natural deduction;
    \item understand sequent calculus;
    \item interpret sequents;
    \item understand structural rules;
    \item recognize weakening, contraction, and exchange;
    \item understand logical introduction rules;
    \item distinguish left and right sequent rules;
    \item understand the cut rule;
    \item recognize cut elimination;
    \item understand the subformula property;
    \item recognize normalization in natural deduction;
    \item distinguish normal and nonnormal proofs;
    \item understand detours in proofs;
    \item recognize proof reduction;
    \item understand consistency proof strategies;
    \item distinguish syntactic and semantic consistency;
    \item understand proof-theoretic strength;
    \item recognize conservative extensions;
    \item understand relative consistency;
    \item understand formalized arithmetic;
    \item recognize induction schemas;
    \item understand finitistic reasoning conceptually;
    \item recognize Gödel coding;
    \item understand arithmetization of syntax;
    \item recognize provability predicates;
    \item understand the diagonal lemma conceptually;
    \item distinguish completeness from incompleteness;
    \item understand Gödel's first incompleteness theorem conceptually;
    \item understand Gödel's second incompleteness theorem conceptually;
    \item recognize the role of consistency assumptions;
    \item understand derivability conditions conceptually;
    \item recognize reflection principles;
    \item understand ordinal analysis conceptually;
    \item recognize proof-theoretic ordinals;
    \item understand the role of \(\varepsilon_0\);
    \item recognize Gentzen-style consistency proofs conceptually;
    \item understand constructive proof theory;
    \item recognize intuitionistic logic;
    \item distinguish classical and intuitionistic principles;
    \item understand the law of excluded middle;
    \item understand double-negation elimination;
    \item recognize Curry--Howard correspondence conceptually;
    \item understand proofs as programs conceptually;
    \item recognize normalization as computation;
    \item understand proof search;
    \item recognize automated deduction;
    \item understand proof complexity conceptually;
    \item distinguish proof length from computational complexity;
    \item recognize resolution proof systems;
    \item understand tableaux conceptually;
    \item recognize formal proof verification;
    \item connect proof theory with logic, computability, type theory,
          constructive mathematics, automated theorem proving, and the
          foundations of mathematics.
\end{itemize}
}

%%%%%%%%%%%%%%%%%%%%%%%%%%%%%%%%%%%%%%%%%%%%%%%%%%%%%%%%%%%%
\subsection{What Is Proof Theory?}
\label{subsec:what-is-proof-theory}
%%%%%%%%%%%%%%%%%%%%%%%%%%%%%%%%%%%%%%%%%%%%%%%%%%%%%%%%%%%%

Proof theory studies formal proofs as mathematical objects.

Instead of asking only

\[
\boxed{
\text{Is a statement true?}
}
\]

proof theory asks questions such as

\[
\boxed{
\text{Can the statement be formally derived?}
}
\]

\[
\boxed{
\text{What rules are required?}
}
\]

\[
\boxed{
\text{Can the proof be simplified?}
}
\]

and

\[
\boxed{
\text{How strong is the proof system?}
}
\]

Thus proof theory studies the syntactic side of mathematical logic.

%%%%%%%%%%%%%%%%%%%%%%%%%%%%%%%%%%%%%%%%%%%%%%%%%%%%%%%%%%%%
\subsection{Proofs as Finite Objects}
\label{subsec:proofs-as-finite-objects}
%%%%%%%%%%%%%%%%%%%%%%%%%%%%%%%%%%%%%%%%%%%%%%%%%%%%%%%%%%%%

A formal proof is typically a finite sequence or finite tree of formulas.

Each line or node is justified by:

\[
\boxed{
\text{an axiom},
}
\]

\[
\boxed{
\text{an assumption},
}
\]

or

\[
\boxed{
\text{an inference rule}.
}
\]

Because proofs are finite symbolic objects, they can themselves be encoded,
analyzed, transformed, and verified.

%%%%%%%%%%%%%%%%%%%%%%%%%%%%%%%%%%%%%%%%%%%%%%%%%%%%%%%%%%%%
\subsection{Object Theory and Metatheory}
\label{subsec:proof-theory-object-metatheory}
%%%%%%%%%%%%%%%%%%%%%%%%%%%%%%%%%%%%%%%%%%%%%%%%%%%%%%%%%%%%

The \textbf{object theory} is the formal system being studied.

The \textbf{metatheory} is the mathematical framework used to reason about
that formal system.

For example,

\[
T\vdash\varphi
\]

is a statement about derivability inside theory \(T\).

A proof that

\[
T
\text{ is consistent}
\]

is generally a metatheoretic argument about \(T\).

%%%%%%%%%%%%%%%%%%%%%%%%%%%%%%%%%%%%%%%%%%%%%%%%%%%%%%%%%%%%
\subsection{Axioms}
\label{subsec:proof-theoretic-axioms}
%%%%%%%%%%%%%%%%%%%%%%%%%%%%%%%%%%%%%%%%%%%%%%%%%%%%%%%%%%%%

An axiom is a formula accepted without derivation inside a formal system.

A theory may contain:

\[
\boxed{
\text{logical axioms}
}
\]

and

\[
\boxed{
\text{nonlogical axioms}.
}
\]

Logical axioms govern logical reasoning.

Nonlogical axioms describe the mathematical subject, such as arithmetic,
set theory, or geometry.

%%%%%%%%%%%%%%%%%%%%%%%%%%%%%%%%%%%%%%%%%%%%%%%%%%%%%%%%%%%%
\subsection{Inference Rules}
\label{subsec:proof-theoretic-inference-rules}
%%%%%%%%%%%%%%%%%%%%%%%%%%%%%%%%%%%%%%%%%%%%%%%%%%%%%%%%%%%%

An inference rule specifies how new formulas may be derived from previous
ones.

For example, modus ponens is

\[
\boxed{
\frac{
\varphi
\qquad
\varphi\rightarrow\psi
}{
\psi
}.
}
\]

The formulas above the line are premises.

The formula below the line is the conclusion.

%%%%%%%%%%%%%%%%%%%%%%%%%%%%%%%%%%%%%%%%%%%%%%%%%%%%%%%%%%%%
\subsection{Derivations}
\label{subsec:formal-derivations-proof-theory}
%%%%%%%%%%%%%%%%%%%%%%%%%%%%%%%%%%%%%%%%%%%%%%%%%%%%%%%%%%%%

A derivation from assumptions \(\Gamma\) is a finite construction ending in
a formula \(\varphi\).

If such a derivation exists, write

\[
\boxed{
\Gamma\vdash\varphi.
}
\]

If no assumptions are used,

\[
\boxed{
\vdash\varphi.
}
\]

The formula is then a theorem of the proof system.

%%%%%%%%%%%%%%%%%%%%%%%%%%%%%%%%%%%%%%%%%%%%%%%%%%%%%%%%%%%%
\subsection{Proof Trees}
\label{subsec:proof-trees}
%%%%%%%%%%%%%%%%%%%%%%%%%%%%%%%%%%%%%%%%%%%%%%%%%%%%%%%%%%%%

Instead of writing proofs as linear sequences, one may display inference
structure as a tree.

For example,

\[
\frac{
P
\qquad
P\rightarrow Q
}{
Q
}
\]

shows immediately that \(Q\) depends on two premises.

More complicated proofs form branching trees.

Proof trees make dependencies explicit.

%%%%%%%%%%%%%%%%%%%%%%%%%%%%%%%%%%%%%%%%%%%%%%%%%%%%%%%%%%%%
\subsection{Hilbert-Style Systems}
\label{subsec:hilbert-style-systems}
%%%%%%%%%%%%%%%%%%%%%%%%%%%%%%%%%%%%%%%%%%%%%%%%%%%%%%%%%%%%

A Hilbert-style system typically contains:

\[
\boxed{
\text{many axiom schemas}
}
\]

and

\[
\boxed{
\text{few inference rules}.
}
\]

Often modus ponens is the main propositional inference rule.

Such systems are compact and mathematically elegant, but formal proofs can
be lengthy and unlike ordinary mathematical reasoning.

%%%%%%%%%%%%%%%%%%%%%%%%%%%%%%%%%%%%%%%%%%%%%%%%%%%%%%%%%%%%
\subsection{Axiom Schemas}
\label{subsec:proof-theory-axiom-schemas}
%%%%%%%%%%%%%%%%%%%%%%%%%%%%%%%%%%%%%%%%%%%%%%%%%%%%%%%%%%%%

An axiom schema represents infinitely many axioms.

For example, the pattern

\[
\boxed{
\varphi
\rightarrow
(\psi\rightarrow\varphi)
}
\]

represents one axiom for every pair of formulas \(\varphi,\psi\).

Thus a finite syntactic pattern can specify infinitely many formal axioms.

%%%%%%%%%%%%%%%%%%%%%%%%%%%%%%%%%%%%%%%%%%%%%%%%%%%%%%%%%%%%
\subsection{Natural Deduction}
\label{subsec:proof-theory-natural-deduction}
%%%%%%%%%%%%%%%%%%%%%%%%%%%%%%%%%%%%%%%%%%%%%%%%%%%%%%%%%%%%

Natural deduction uses rules designed to resemble ordinary mathematical
reasoning.

Each connective has:

\[
\boxed{
\text{introduction rules}
}
\]

and

\[
\boxed{
\text{elimination rules}.
}
\]

Introduction rules describe how to prove a formula containing a connective.

Elimination rules describe how to use such a formula.

%%%%%%%%%%%%%%%%%%%%%%%%%%%%%%%%%%%%%%%%%%%%%%%%%%%%%%%%%%%%
\subsection{Conjunction Rules}
\label{subsec:natural-deduction-conjunction}
%%%%%%%%%%%%%%%%%%%%%%%%%%%%%%%%%%%%%%%%%%%%%%%%%%%%%%%%%%%%

Conjunction introduction is

\[
\boxed{
\frac{
\varphi
\qquad
\psi
}{
\varphi\land\psi
}.
}
\]

Conjunction elimination gives

\[
\boxed{
\frac{
\varphi\land\psi
}{
\varphi
}
}
\]

and

\[
\boxed{
\frac{
\varphi\land\psi
}{
\psi
}.
}
\]

%%%%%%%%%%%%%%%%%%%%%%%%%%%%%%%%%%%%%%%%%%%%%%%%%%%%%%%%%%%%
\subsection{Implication Rules}
\label{subsec:natural-deduction-implication}
%%%%%%%%%%%%%%%%%%%%%%%%%%%%%%%%%%%%%%%%%%%%%%%%%%%%%%%%%%%%

Implication elimination is modus ponens:

\[
\boxed{
\frac{
\varphi
\qquad
\varphi\rightarrow\psi
}{
\psi
}.
}
\]

Implication introduction temporarily assumes \(\varphi\), derives \(\psi\),
and then discharges the assumption:

\[
\boxed{
\frac{
[\varphi]
\quad
\vdots
\quad
\psi
}{
\varphi\rightarrow\psi
}.
}
\]

%%%%%%%%%%%%%%%%%%%%%%%%%%%%%%%%%%%%%%%%%%%%%%%%%%%%%%%%%%%%
\subsection{Worked Example: Natural Deduction}
\label{subsec:worked-natural-deduction}
%%%%%%%%%%%%%%%%%%%%%%%%%%%%%%%%%%%%%%%%%%%%%%%%%%%%%%%%%%%%

\begin{workedexamplebox}{Deriving \(P\rightarrow(Q\rightarrow P)\)}

Assume

\[
P.
\]

Now assume

\[
Q.
\]

From the first assumption, we already have

\[
P.
\]

Discharge the assumption \(Q\):

\[
Q\rightarrow P.
\]

Then discharge the assumption \(P\):

\[
P\rightarrow(Q\rightarrow P).
\]

\begin{answerbox}
\[
\boxed{
\vdash
P\rightarrow(Q\rightarrow P).
}
\]
\end{answerbox}

\end{workedexamplebox}

%%%%%%%%%%%%%%%%%%%%%%%%%%%%%%%%%%%%%%%%%%%%%%%%%%%%%%%%%%%%
\subsection{Disjunction Rules}
\label{subsec:natural-deduction-disjunction}
%%%%%%%%%%%%%%%%%%%%%%%%%%%%%%%%%%%%%%%%%%%%%%%%%%%%%%%%%%%%

Disjunction introduction gives

\[
\boxed{
\frac{
\varphi
}{
\varphi\lor\psi
}
}
\]

and

\[
\boxed{
\frac{
\psi
}{
\varphi\lor\psi
}.
}
\]

Disjunction elimination formalizes proof by cases.

If

\[
\varphi\lor\psi
\]

holds, and both \(\varphi\) and \(\psi\) separately imply \(\chi\), then
\(\chi\) follows.

%%%%%%%%%%%%%%%%%%%%%%%%%%%%%%%%%%%%%%%%%%%%%%%%%%%%%%%%%%%%
\subsection{Negation Rules}
\label{subsec:natural-deduction-negation}
%%%%%%%%%%%%%%%%%%%%%%%%%%%%%%%%%%%%%%%%%%%%%%%%%%%%%%%%%%%%

Negation can be represented using falsity:

\[
\boxed{
\neg\varphi
\equiv
\varphi\rightarrow\bot.
}
\]

If assuming \(\varphi\) leads to contradiction, derive

\[
\boxed{
\neg\varphi.
}
\]

This formalizes proof of a negative statement.

%%%%%%%%%%%%%%%%%%%%%%%%%%%%%%%%%%%%%%%%%%%%%%%%%%%%%%%%%%%%
\subsection{Classical Reasoning}
\label{subsec:classical-proof-reasoning}
%%%%%%%%%%%%%%%%%%%%%%%%%%%%%%%%%%%%%%%%%%%%%%%%%%%%%%%%%%%%

Classical logic permits principles such as:

\[
\boxed{
\varphi\lor\neg\varphi
}
\]

and

\[
\boxed{
\neg\neg\varphi
\rightarrow
\varphi.
}
\]

These are the law of excluded middle and double-negation elimination.

They are not generally accepted as unrestricted principles in
intuitionistic logic.

%%%%%%%%%%%%%%%%%%%%%%%%%%%%%%%%%%%%%%%%%%%%%%%%%%%%%%%%%%%%
\subsection{Sequent Calculus}
\label{subsec:sequent-calculus}
%%%%%%%%%%%%%%%%%%%%%%%%%%%%%%%%%%%%%%%%%%%%%%%%%%%%%%%%%%%%

Sequent calculus makes the structure of assumptions and conclusions
explicit.

A sequent has form

\[
\boxed{
\Gamma
\Rightarrow
\Delta,
}
\]

where \(\Gamma\) and \(\Delta\) are collections of formulas.

Intuitively, the sequent states that the formulas on the left entail those
on the right.

%%%%%%%%%%%%%%%%%%%%%%%%%%%%%%%%%%%%%%%%%%%%%%%%%%%%%%%%%%%%
\subsection{Meaning of a Sequent}
\label{subsec:meaning-of-sequent}
%%%%%%%%%%%%%%%%%%%%%%%%%%%%%%%%%%%%%%%%%%%%%%%%%%%%%%%%%%%%

In classical logic,

\[
\Gamma
\Rightarrow
\Delta
\]

can be interpreted semantically as

\[
\boxed{
\bigwedge_{\varphi\in\Gamma}\varphi
\rightarrow
\bigvee_{\psi\in\Delta}\psi.
}
\]

Thus the left side behaves conjunctively and the right side
disjunctively.

%%%%%%%%%%%%%%%%%%%%%%%%%%%%%%%%%%%%%%%%%%%%%%%%%%%%%%%%%%%%
\subsection{Initial Sequents}
\label{subsec:initial-sequents}
%%%%%%%%%%%%%%%%%%%%%%%%%%%%%%%%%%%%%%%%%%%%%%%%%%%%%%%%%%%%

A basic identity sequent is

\[
\boxed{
A
\Rightarrow
A.
}
\]

This expresses the trivial fact that assumption \(A\) entails \(A\).

More complex proofs are generated by applying sequent rules.

%%%%%%%%%%%%%%%%%%%%%%%%%%%%%%%%%%%%%%%%%%%%%%%%%%%%%%%%%%%%
\subsection{Structural Rules}
\label{subsec:structural-rules}
%%%%%%%%%%%%%%%%%%%%%%%%%%%%%%%%%%%%%%%%%%%%%%%%%%%%%%%%%%%%

Structural rules manipulate the arrangement of assumptions or conclusions
without referring to their internal logical form.

The principal structural rules are:

\[
\boxed{
\text{weakening},
}
\]

\[
\boxed{
\text{contraction},
}
\]

and

\[
\boxed{
\text{exchange}.
}
\]

Different logics may permit or restrict these rules.

%%%%%%%%%%%%%%%%%%%%%%%%%%%%%%%%%%%%%%%%%%%%%%%%%%%%%%%%%%%%
\subsection{Weakening}
\label{subsec:weakening}
%%%%%%%%%%%%%%%%%%%%%%%%%%%%%%%%%%%%%%%%%%%%%%%%%%%%%%%%%%%%

Left weakening has schematic form

\[
\boxed{
\frac{
\Gamma\Rightarrow\Delta
}{
\Gamma,A\Rightarrow\Delta
}.
}
\]

It says that adding an unused assumption does not invalidate a derivation.

Similarly, right weakening permits additional conclusions in classical
multi-conclusion systems.

%%%%%%%%%%%%%%%%%%%%%%%%%%%%%%%%%%%%%%%%%%%%%%%%%%%%%%%%%%%%
\subsection{Contraction}
\label{subsec:contraction}
%%%%%%%%%%%%%%%%%%%%%%%%%%%%%%%%%%%%%%%%%%%%%%%%%%%%%%%%%%%%

Contraction removes duplicate occurrences:

\[
\boxed{
\frac{
\Gamma,A,A\Rightarrow\Delta
}{
\Gamma,A\Rightarrow\Delta
}.
}
\]

This means an assumption may be reused without being treated as a
consumable resource.

Linear logic famously modifies this structural behavior.

%%%%%%%%%%%%%%%%%%%%%%%%%%%%%%%%%%%%%%%%%%%%%%%%%%%%%%%%%%%%
\subsection{Exchange}
\label{subsec:exchange}
%%%%%%%%%%%%%%%%%%%%%%%%%%%%%%%%%%%%%%%%%%%%%%%%%%%%%%%%%%%%

Exchange permits reordering formulas:

\[
\boxed{
\Gamma,A,B,\Pi
\Rightarrow
\Delta
}
\]

can be rearranged to

\[
\boxed{
\Gamma,B,A,\Pi
\Rightarrow
\Delta.
}
\]

In standard classical and intuitionistic systems, order of assumptions does
not matter.

%%%%%%%%%%%%%%%%%%%%%%%%%%%%%%%%%%%%%%%%%%%%%%%%%%%%%%%%%%%%
\subsection{Logical Rules in Sequent Calculus}
\label{subsec:logical-sequent-rules}
%%%%%%%%%%%%%%%%%%%%%%%%%%%%%%%%%%%%%%%%%%%%%%%%%%%%%%%%%%%%

Logical rules introduce connectives on either side of a sequent.

For example, conjunction right introduction has schematic form

\[
\boxed{
\frac{
\Gamma\Rightarrow\Delta,A
\qquad
\Gamma\Rightarrow\Delta,B
}{
\Gamma\Rightarrow\Delta,A\land B
}.
}
\]

The exact rule presentation depends on the sequent system.

%%%%%%%%%%%%%%%%%%%%%%%%%%%%%%%%%%%%%%%%%%%%%%%%%%%%%%%%%%%%
\subsection{Left and Right Rules}
\label{subsec:left-right-sequent-rules}
%%%%%%%%%%%%%%%%%%%%%%%%%%%%%%%%%%%%%%%%%%%%%%%%%%%%%%%%%%%%

A left rule explains how a logical connective behaves when it appears among
the assumptions.

A right rule explains how it behaves among the conclusions.

Thus sequent calculus exposes a symmetry between:

\[
\boxed{
\text{using a formula}
}
\]

and

\[
\boxed{
\text{proving a formula}.
}
\]

%%%%%%%%%%%%%%%%%%%%%%%%%%%%%%%%%%%%%%%%%%%%%%%%%%%%%%%%%%%%
\subsection{The Cut Rule}
\label{subsec:cut-rule}
%%%%%%%%%%%%%%%%%%%%%%%%%%%%%%%%%%%%%%%%%%%%%%%%%%%%%%%%%%%%

The cut rule has form

\[
\boxed{
\frac{
\Gamma\Rightarrow\Delta,A
\qquad
A,\Pi\Rightarrow\Lambda
}{
\Gamma,\Pi
\Rightarrow
\Delta,\Lambda
}.
}
\]

The intermediate formula \(A\) disappears from the conclusion.

Conceptually, if one proof derives \(A\) and another uses \(A\), the two
proofs can be composed.

%%%%%%%%%%%%%%%%%%%%%%%%%%%%%%%%%%%%%%%%%%%%%%%%%%%%%%%%%%%%
\subsection{Cut as a Lemma}
\label{subsec:cut-as-lemma}
%%%%%%%%%%%%%%%%%%%%%%%%%%%%%%%%%%%%%%%%%%%%%%%%%%%%%%%%%%%%

The cut formula acts like an intermediate lemma.

One first proves

\[
A,
\]

then uses

\[
A
\]

to prove the final result.

Thus cut formalizes ordinary mathematical reasoning with lemmas.

Surprisingly, in standard sequent calculi, cut can often be eliminated.

%%%%%%%%%%%%%%%%%%%%%%%%%%%%%%%%%%%%%%%%%%%%%%%%%%%%%%%%%%%%
\subsection{Cut Elimination}
\label{subsec:cut-elimination}
%%%%%%%%%%%%%%%%%%%%%%%%%%%%%%%%%%%%%%%%%%%%%%%%%%%%%%%%%%%%

\begin{theorem}
In standard classical and intuitionistic first-order sequent calculi, every
proof using the cut rule can be transformed into a proof of the same sequent
without cut.
\end{theorem}

This is Gentzen's cut-elimination theorem.

Symbolically,

\[
\boxed{
\text{proof with cuts}
\longrightarrow
\text{cut-free proof}.
}
\]

%%%%%%%%%%%%%%%%%%%%%%%%%%%%%%%%%%%%%%%%%%%%%%%%%%%%%%%%%%%%
\subsection{Why Cut Elimination Matters}
\label{subsec:why-cut-elimination-matters}
%%%%%%%%%%%%%%%%%%%%%%%%%%%%%%%%%%%%%%%%%%%%%%%%%%%%%%%%%%%%

Cut elimination has several consequences.

It reveals that intermediate formulas are not logically necessary for
derivability.

It also supports:

\[
\boxed{
\text{consistency proofs},
}
\]

\[
\boxed{
\text{subformula properties},
}
\]

and

\[
\boxed{
\text{proof search}.
}
\]

It is one of the central structural results of proof theory.

%%%%%%%%%%%%%%%%%%%%%%%%%%%%%%%%%%%%%%%%%%%%%%%%%%%%%%%%%%%%
\subsection{Subformula Property}
\label{subsec:subformula-property}
%%%%%%%%%%%%%%%%%%%%%%%%%%%%%%%%%%%%%%%%%%%%%%%%%%%%%%%%%%%%

In many cut-free proof systems, every formula occurring in a proof is built
from subformulas of the final sequent.

Thus the proof does not introduce arbitrary unrelated intermediate
statements.

Conceptually,

\[
\boxed{
\text{cut-free proof}
\Longrightarrow
\text{strong control of formula complexity}.
}
\]

%%%%%%%%%%%%%%%%%%%%%%%%%%%%%%%%%%%%%%%%%%%%%%%%%%%%%%%%%%%%
\subsection{Proof Search}
\label{subsec:proof-search}
%%%%%%%%%%%%%%%%%%%%%%%%%%%%%%%%%%%%%%%%%%%%%%%%%%%%%%%%%%%%

The subformula property makes proof search more manageable.

Instead of guessing arbitrary lemmas, a proof-search procedure can often
restrict attention to formulas already present in the goal.

Thus sequent calculus provides an algorithmic view of logical deduction.

%%%%%%%%%%%%%%%%%%%%%%%%%%%%%%%%%%%%%%%%%%%%%%%%%%%%%%%%%%%%
\subsection{Normalization}
\label{subsec:proof-normalization}
%%%%%%%%%%%%%%%%%%%%%%%%%%%%%%%%%%%%%%%%%%%%%%%%%%%%%%%%%%%%

Natural-deduction proofs may contain unnecessary detours.

For example, one may introduce a connective and immediately eliminate it.

Normalization transforms proofs to remove such detours.

Conceptually,

\[
\boxed{
\text{nonnormal proof}
\longrightarrow
\text{normal proof}.
}
\]

%%%%%%%%%%%%%%%%%%%%%%%%%%%%%%%%%%%%%%%%%%%%%%%%%%%%%%%%%%%%
\subsection{Introduction-Elimination Detours}
\label{subsec:introduction-elimination-detours}
%%%%%%%%%%%%%%%%%%%%%%%%%%%%%%%%%%%%%%%%%%%%%%%%%%%%%%%%%%%%

Suppose one derives

\[
A\land B
\]

using conjunction introduction and immediately extracts \(A\).

Schematically,

\[
\frac{
A
\qquad
B
}{
A\land B
}
\]

followed by

\[
\frac{
A\land B
}{
A
}.
\]

The intermediate conjunction is unnecessary.

Normalization removes this detour.

%%%%%%%%%%%%%%%%%%%%%%%%%%%%%%%%%%%%%%%%%%%%%%%%%%%%%%%%%%%%
\subsection{Worked Example: Proof Reduction}
\label{subsec:worked-proof-reduction}
%%%%%%%%%%%%%%%%%%%%%%%%%%%%%%%%%%%%%%%%%%%%%%%%%%%%%%%%%%%%

\begin{workedexamplebox}{Removing a Conjunction Detour}

Suppose a proof already contains derivations of

\[
A
\]

and

\[
B.
\]

One constructs

\[
A\land B
\]

and immediately applies conjunction elimination to recover \(A\).

The entire two-step sequence can be replaced by the original derivation of
\(A\).

\begin{answerbox}
\[
\boxed{
A
\quad\text{is obtained directly without introducing }A\land B.
}
\]
\end{answerbox}

This is a simple normalization step.

\end{workedexamplebox}

%%%%%%%%%%%%%%%%%%%%%%%%%%%%%%%%%%%%%%%%%%%%%%%%%%%%%%%%%%%%
\subsection{Normal Forms of Proofs}
\label{subsec:normal-forms-proofs}
%%%%%%%%%%%%%%%%%%%%%%%%%%%%%%%%%%%%%%%%%%%%%%%%%%%%%%%%%%%%

A normal proof contains no reducible introduction-elimination detours.

Normalization provides a canonical or simplified structural form for
derivations.

The analogy with computation is

\[
\boxed{
\text{proof reduction}
\leftrightarrow
\text{program evaluation}.
}
\]

This connection becomes central in type theory.

%%%%%%%%%%%%%%%%%%%%%%%%%%%%%%%%%%%%%%%%%%%%%%%%%%%%%%%%%%%%
\subsection{Consistency}
\label{subsec:proof-theoretic-consistency}
%%%%%%%%%%%%%%%%%%%%%%%%%%%%%%%%%%%%%%%%%%%%%%%%%%%%%%%%%%%%

A formal theory \(T\) is consistent if contradiction is not derivable:

\[
\boxed{
T\nvdash\bot.
}
\]

In standard classical systems, equivalently, there is no formula
\(\varphi\) such that

\[
T\vdash\varphi
\]

and

\[
T\vdash\neg\varphi.
\]

%%%%%%%%%%%%%%%%%%%%%%%%%%%%%%%%%%%%%%%%%%%%%%%%%%%%%%%%%%%%
\subsection{Syntactic Versus Semantic Consistency}
\label{subsec:syntactic-semantic-consistency}
%%%%%%%%%%%%%%%%%%%%%%%%%%%%%%%%%%%%%%%%%%%%%%%%%%%%%%%%%%%%

Syntactic consistency means

\[
\boxed{
T\nvdash\bot.
}
\]

Semantic satisfiability means

\[
\boxed{
\exists\mathcal{M}
\,
(\mathcal{M}\models T).
}
\]

For standard first-order logic, soundness and completeness connect these
notions:

\[
\boxed{
T
\text{ consistent}
\iff
T
\text{ satisfiable}.
}
\]

%%%%%%%%%%%%%%%%%%%%%%%%%%%%%%%%%%%%%%%%%%%%%%%%%%%%%%%%%%%%
\subsection{Proof-Theoretic Consistency Proofs}
\label{subsec:proof-theoretic-consistency-proofs}
%%%%%%%%%%%%%%%%%%%%%%%%%%%%%%%%%%%%%%%%%%%%%%%%%%%%%%%%%%%%

A proof-theoretic consistency argument analyzes possible formal derivations
and shows that no derivation of

\[
\bot
\]

can exist.

Typical tools include:

\[
\boxed{
\text{normalization},
}
\]

\[
\boxed{
\text{cut elimination},
}
\]

and

\[
\boxed{
\text{ordinal measures}.
}
\]

%%%%%%%%%%%%%%%%%%%%%%%%%%%%%%%%%%%%%%%%%%%%%%%%%%%%%%%%%%%%
\subsection{Consistency From Cut Elimination}
\label{subsec:consistency-cut-elimination}
%%%%%%%%%%%%%%%%%%%%%%%%%%%%%%%%%%%%%%%%%%%%%%%%%%%%%%%%%%%%

Suppose a sequent system admits cut elimination.

If the empty contradiction sequent had a proof, it would have a cut-free
proof.

The structure of cut-free proofs may then show that no final inference can
produce such a sequent.

Thus

\[
\boxed{
\text{cut elimination}
\Longrightarrow
\text{consistency}
}
\]

for appropriate calculi.

%%%%%%%%%%%%%%%%%%%%%%%%%%%%%%%%%%%%%%%%%%%%%%%%%%%%%%%%%%%%
\subsection{Conservative Extensions}
\label{subsec:conservative-extensions}
%%%%%%%%%%%%%%%%%%%%%%%%%%%%%%%%%%%%%%%%%%%%%%%%%%%%%%%%%%%%

Suppose

\[
T\subseteq T'.
\]

The stronger theory \(T'\) is conservative over \(T\) for a class of
sentences if whenever

\[
T'\vdash\varphi
\]

for a sentence \(\varphi\) in that class, already

\[
\boxed{
T\vdash\varphi.
}
\]

Thus the extension adds expressive or proof-theoretic convenience without
creating new theorems of the specified old kind.

%%%%%%%%%%%%%%%%%%%%%%%%%%%%%%%%%%%%%%%%%%%%%%%%%%%%%%%%%%%%
\subsection{Relative Consistency}
\label{subsec:proof-theoretic-relative-consistency}
%%%%%%%%%%%%%%%%%%%%%%%%%%%%%%%%%%%%%%%%%%%%%%%%%%%%%%%%%%%%

A relative consistency proof establishes a statement such as

\[
\boxed{
\operatorname{Con}(T)
\Longrightarrow
\operatorname{Con}(T').
}
\]

This says that if the base theory \(T\) is free of contradiction, then so
is \(T'\).

Relative consistency does not by itself establish absolute consistency.

%%%%%%%%%%%%%%%%%%%%%%%%%%%%%%%%%%%%%%%%%%%%%%%%%%%%%%%%%%%%
\subsection{Proof-Theoretic Strength}
\label{subsec:proof-theoretic-strength}
%%%%%%%%%%%%%%%%%%%%%%%%%%%%%%%%%%%%%%%%%%%%%%%%%%%%%%%%%%%%

Two formal theories may differ in which statements they can prove.

A stronger theory may prove:

\[
\boxed{
\text{more induction principles},
}
\]

\[
\boxed{
\text{more transfinite recursion},
}
\]

or

\[
\boxed{
\text{consistency statements for weaker theories}.
}
\]

Proof theory seeks ways to compare these strengths systematically.

%%%%%%%%%%%%%%%%%%%%%%%%%%%%%%%%%%%%%%%%%%%%%%%%%%%%%%%%%%%%
\subsection{Formal Arithmetic}
\label{subsec:formal-arithmetic-proof-theory}
%%%%%%%%%%%%%%%%%%%%%%%%%%%%%%%%%%%%%%%%%%%%%%%%%%%%%%%%%%%%

Arithmetic can be formalized in a first-order language containing symbols
such as

\[
\boxed{
0,
\qquad
S,
\qquad
+,
\qquad
\cdot.
}
\]

Here \(S\) denotes successor.

A formal arithmetic theory contains axioms describing these operations and
an induction principle.

%%%%%%%%%%%%%%%%%%%%%%%%%%%%%%%%%%%%%%%%%%%%%%%%%%%%%%%%%%%%
\subsection{Induction Schema}
\label{subsec:induction-schema-proof-theory}
%%%%%%%%%%%%%%%%%%%%%%%%%%%%%%%%%%%%%%%%%%%%%%%%%%%%%%%%%%%%

For formula \(\varphi(x)\), induction has form

\[
\boxed{
\left[
\varphi(0)
\land
\forall n
(
\varphi(n)\rightarrow\varphi(Sn)
)
\right]
\rightarrow
\forall n\,\varphi(n).
}
\]

Because one such axiom is included for each appropriate formula
\(\varphi\), induction is an axiom schema in first-order arithmetic.

%%%%%%%%%%%%%%%%%%%%%%%%%%%%%%%%%%%%%%%%%%%%%%%%%%%%%%%%%%%%
\subsection{Restricted Induction}
\label{subsec:restricted-induction}
%%%%%%%%%%%%%%%%%%%%%%%%%%%%%%%%%%%%%%%%%%%%%%%%%%%%%%%%%%%%

One can form weaker arithmetic theories by allowing induction only for
restricted classes of formulas.

For example, one may restrict induction according to logical complexity.

Such theories help measure how much mathematical reasoning depends on
particular induction principles.

%%%%%%%%%%%%%%%%%%%%%%%%%%%%%%%%%%%%%%%%%%%%%%%%%%%%%%%%%%%%
\subsection{Gödel Coding}
\label{subsec:godel-coding}
%%%%%%%%%%%%%%%%%%%%%%%%%%%%%%%%%%%%%%%%%%%%%%%%%%%%%%%%%%%%

Formal formulas and proofs are finite symbolic objects.

Gödel's key idea is that these objects can be encoded by natural numbers.

Conceptually,

\[
\boxed{
\text{symbol string}
\longrightarrow
\text{natural-number code}.
}
\]

Thus arithmetic can reason indirectly about syntax.

%%%%%%%%%%%%%%%%%%%%%%%%%%%%%%%%%%%%%%%%%%%%%%%%%%%%%%%%%%%%
\subsection{Arithmetization of Syntax}
\label{subsec:arithmetization-of-syntax}
%%%%%%%%%%%%%%%%%%%%%%%%%%%%%%%%%%%%%%%%%%%%%%%%%%%%%%%%%%%%

Once formulas and proofs are encoded numerically, syntactic relations can
be represented arithmetically.

Examples include:

\[
\boxed{
\operatorname{Formula}(n),
}
\]

meaning \(n\) codes a formula, and

\[
\boxed{
\operatorname{Proof}_T(p,n),
}
\]

meaning \(p\) codes a \(T\)-proof of the formula coded by \(n\).

This is called arithmetization of syntax.

%%%%%%%%%%%%%%%%%%%%%%%%%%%%%%%%%%%%%%%%%%%%%%%%%%%%%%%%%%%%
\subsection{Provability Predicate}
\label{subsec:provability-predicate}
%%%%%%%%%%%%%%%%%%%%%%%%%%%%%%%%%%%%%%%%%%%%%%%%%%%%%%%%%%%%

Define conceptually

\[
\boxed{
\operatorname{Pr}_T(n)
}
\]

to mean that there exists a code \(p\) for a \(T\)-proof of the formula
with code \(n\).

Schematically,

\[
\boxed{
\operatorname{Pr}_T(n)
\equiv
\exists p\,
\operatorname{Proof}_T(p,n).
}
\]

Thus formal provability becomes an arithmetic property.

%%%%%%%%%%%%%%%%%%%%%%%%%%%%%%%%%%%%%%%%%%%%%%%%%%%%%%%%%%%%
\subsection{Self-Reference}
\label{subsec:logical-self-reference}
%%%%%%%%%%%%%%%%%%%%%%%%%%%%%%%%%%%%%%%%%%%%%%%%%%%%%%%%%%%%

Once syntax is encoded inside arithmetic, formulas can indirectly refer to
their own codes.

This makes rigorous forms of self-reference possible.

The diagonal lemma formalizes this mechanism.

%%%%%%%%%%%%%%%%%%%%%%%%%%%%%%%%%%%%%%%%%%%%%%%%%%%%%%%%%%%%
\subsection{Diagonal Lemma}
\label{subsec:diagonal-lemma}
%%%%%%%%%%%%%%%%%%%%%%%%%%%%%%%%%%%%%%%%%%%%%%%%%%%%%%%%%%%%

Conceptually, the diagonal lemma states that for a sufficiently expressive
theory and formula

\[
\varphi(x),
\]

there exists a sentence \(G\) satisfying

\[
\boxed{
T\vdash
G
\leftrightarrow
\varphi(\ulcorner G\urcorner),
}
\]

where

\[
\ulcorner G\urcorner
\]

denotes the Gödel code of \(G\).

This creates a formal sentence that speaks about itself indirectly.

%%%%%%%%%%%%%%%%%%%%%%%%%%%%%%%%%%%%%%%%%%%%%%%%%%%%%%%%%%%%
\subsection{Gödel's First Incompleteness Theorem}
\label{subsec:first-incompleteness}
%%%%%%%%%%%%%%%%%%%%%%%%%%%%%%%%%%%%%%%%%%%%%%%%%%%%%%%%%%%%

Very roughly, Gödel's First Incompleteness Theorem states that any
consistent, effectively axiomatized formal theory sufficiently strong to
represent elementary arithmetic cannot decide every arithmetic sentence.

Under standard hypotheses, there exists a sentence \(G\) such that

\[
\boxed{
T\nvdash G
}
\]

and, with appropriate strengthening of consistency assumptions in some
classical formulations,

\[
\boxed{
T\nvdash\neg G.
}
\]

Modern formulations can sharpen the required assumptions.

%%%%%%%%%%%%%%%%%%%%%%%%%%%%%%%%%%%%%%%%%%%%%%%%%%%%%%%%%%%%
\subsection{What Incompleteness Does Not Say}
\label{subsec:what-incompleteness-does-not-say}
%%%%%%%%%%%%%%%%%%%%%%%%%%%%%%%%%%%%%%%%%%%%%%%%%%%%%%%%%%%%

Gödel incompleteness does not mean:

\[
\boxed{
\text{all mathematics is uncertain},
}
\]

or

\[
\boxed{
\text{every statement is undecidable}.
}
\]

It says that sufficiently strong effective consistent formal theories
cannot decide every sentence expressible in their arithmetic language.

Many individual statements remain perfectly provable.

%%%%%%%%%%%%%%%%%%%%%%%%%%%%%%%%%%%%%%%%%%%%%%%%%%%%%%%%%%%%
\subsection{Logical Completeness Versus Gödel Incompleteness}
\label{subsec:completeness-vs-incompleteness}
%%%%%%%%%%%%%%%%%%%%%%%%%%%%%%%%%%%%%%%%%%%%%%%%%%%%%%%%%%%%

Gödel's completeness theorem states:

\[
\boxed{
T\models\varphi
\Longrightarrow
T\vdash\varphi
}
\]

for first-order logical consequence.

Gödel's incompleteness theorem says certain particular theories \(T\) do
not decide every sentence:

\[
\boxed{
T\nvdash\varphi
\quad\text{and}\quad
T\nvdash\neg\varphi
}
\]

for some \(\varphi\).

The two results address different notions of completeness.

%%%%%%%%%%%%%%%%%%%%%%%%%%%%%%%%%%%%%%%%%%%%%%%%%%%%%%%%%%%%
\subsection{Worked Example: Two Different Completeness Notions}
\label{subsec:worked-two-completeness-notions}
%%%%%%%%%%%%%%%%%%%%%%%%%%%%%%%%%%%%%%%%%%%%%%%%%%%%%%%%%%%%

\begin{workedexamplebox}{No Contradiction Between the Theorems}

First-order logical completeness says:

if every model of \(T\) satisfies \(\varphi\), then \(\varphi\) has a formal
proof from \(T\).

Theory incompleteness says:

there can exist a sentence \(\varphi\) such that some models of \(T\)
satisfy \(\varphi\) while others satisfy \(\neg\varphi\).

Then neither sentence is a semantic consequence of \(T\).

\begin{answerbox}
\[
\boxed{
\text{Logical completeness and theory incompleteness are compatible.}
}
\]
\end{answerbox}

\end{workedexamplebox}

%%%%%%%%%%%%%%%%%%%%%%%%%%%%%%%%%%%%%%%%%%%%%%%%%%%%%%%%%%%%
\subsection{Gödel's Second Incompleteness Theorem}
\label{subsec:second-incompleteness}
%%%%%%%%%%%%%%%%%%%%%%%%%%%%%%%%%%%%%%%%%%%%%%%%%%%%%%%%%%%%

Under standard hypotheses, Gödel's Second Incompleteness Theorem states
that a sufficiently strong consistent effectively axiomatized theory
\(T\) cannot prove its own suitably formalized consistency statement.

Schematically,

\[
\boxed{
T
\text{ consistent}
\Longrightarrow
T\nvdash
\operatorname{Con}(T).
}
\]

The precise statement depends on how proof and consistency are represented
inside \(T\).

%%%%%%%%%%%%%%%%%%%%%%%%%%%%%%%%%%%%%%%%%%%%%%%%%%%%%%%%%%%%
\subsection{Consistency Statement}
\label{subsec:formal-consistency-statement}
%%%%%%%%%%%%%%%%%%%%%%%%%%%%%%%%%%%%%%%%%%%%%%%%%%%%%%%%%%%%

A formal consistency statement can be expressed as

\[
\boxed{
\operatorname{Con}(T)
\equiv
\neg
\operatorname{Pr}_T
(
\ulcorner\bot\urcorner
).
}
\]

It states that there is no \(T\)-proof of contradiction.

Because provability is arithmetized, the consistency statement itself can
be represented arithmetically.

%%%%%%%%%%%%%%%%%%%%%%%%%%%%%%%%%%%%%%%%%%%%%%%%%%%%%%%%%%%%
\subsection{Why the Second Incompleteness Theorem Matters}
\label{subsec:why-second-incompleteness-matters}
%%%%%%%%%%%%%%%%%%%%%%%%%%%%%%%%%%%%%%%%%%%%%%%%%%%%%%%%%%%%

The theorem limits certain kinds of self-verification.

A sufficiently strong theory cannot, assuming consistency, establish its
own consistency using only its own internal resources in the standard
formalized sense.

Consistency proofs therefore often require reasoning that is stronger or
different from the theory being analyzed.

%%%%%%%%%%%%%%%%%%%%%%%%%%%%%%%%%%%%%%%%%%%%%%%%%%%%%%%%%%%%
\subsection{Hilbert's Program}
\label{subsec:hilberts-program}
%%%%%%%%%%%%%%%%%%%%%%%%%%%%%%%%%%%%%%%%%%%%%%%%%%%%%%%%%%%%

Hilbert's program sought secure foundations for mathematics by:

\begin{enumerate}
    \item formalizing mathematics;
    \item proving formal systems consistent;
    \item using methods regarded as finitistically reliable.
\end{enumerate}

Gödel's incompleteness theorems showed that the strongest version of this
program cannot be achieved as originally envisioned.

Proof theory nevertheless grew directly from this foundational program.

%%%%%%%%%%%%%%%%%%%%%%%%%%%%%%%%%%%%%%%%%%%%%%%%%%%%%%%%%%%%
\subsection{Finitistic Reasoning}
\label{subsec:finitistic-reasoning}
%%%%%%%%%%%%%%%%%%%%%%%%%%%%%%%%%%%%%%%%%%%%%%%%%%%%%%%%%%%%

Finitistic reasoning emphasizes finite symbolic objects and explicitly
constructible operations.

Formal proofs are particularly suited to this viewpoint because they are
finite strings or trees.

The exact boundary of finitistic mathematics is philosophically and
mathematically subtle.

%%%%%%%%%%%%%%%%%%%%%%%%%%%%%%%%%%%%%%%%%%%%%%%%%%%%%%%%%%%%
\subsection{Gentzen's Consistency Program}
\label{subsec:gentzen-consistency-program}
%%%%%%%%%%%%%%%%%%%%%%%%%%%%%%%%%%%%%%%%%%%%%%%%%%%%%%%%%%%%

Gentzen developed consistency proofs based on structural proof
transformations and transfinite induction.

For arithmetic, the proof assigns ordinal measures to derivations and shows
that proof reduction strictly decreases those measures.

An infinite descent is then impossible.

This connects proof normalization with ordinal theory.

%%%%%%%%%%%%%%%%%%%%%%%%%%%%%%%%%%%%%%%%%%%%%%%%%%%%%%%%%%%%
\subsection{Proof-Theoretic Ordinals}
\label{subsec:proof-theoretic-ordinals}
%%%%%%%%%%%%%%%%%%%%%%%%%%%%%%%%%%%%%%%%%%%%%%%%%%%%%%%%%%%%

A proof-theoretic ordinal is an ordinal used to measure the strength of a
formal theory.

Very roughly, it indicates how much transfinite induction is required to
justify proof reductions associated with the theory.

Different theories correspond to different levels of ordinal strength.

%%%%%%%%%%%%%%%%%%%%%%%%%%%%%%%%%%%%%%%%%%%%%%%%%%%%%%%%%%%%
\subsection{The Ordinal \(\varepsilon_0\)}
\label{subsec:epsilon-zero}
%%%%%%%%%%%%%%%%%%%%%%%%%%%%%%%%%%%%%%%%%%%%%%%%%%%%%%%%%%%%

The ordinal

\[
\boxed{
\varepsilon_0
}
\]

is the least ordinal satisfying

\[
\boxed{
\omega^{\varepsilon_0}
=
\varepsilon_0.
}
\]

It can be viewed as the limit of the sequence

\[
\omega,
\]

\[
\omega^\omega,
\]

\[
\omega^{\omega^\omega},
\]

and continuing iterated exponentiation.

It plays a central role in ordinal analyses related to first-order
arithmetic.

%%%%%%%%%%%%%%%%%%%%%%%%%%%%%%%%%%%%%%%%%%%%%%%%%%%%%%%%%%%%
\subsection{Ordinal Descent in Consistency Proofs}
\label{subsec:ordinal-descent}
%%%%%%%%%%%%%%%%%%%%%%%%%%%%%%%%%%%%%%%%%%%%%%%%%%%%%%%%%%%%

Suppose every hypothetical proof of contradiction is assigned an ordinal

\[
\alpha.
\]

A proof-reduction procedure transforms it into another hypothetical proof
of contradiction with smaller ordinal

\[
\beta<\alpha.
\]

Repeated reduction would produce

\[
\alpha_0
>
\alpha_1
>
\alpha_2
>
\cdots.
\]

Well-foundedness of the ordinals rules out such an infinite descending
sequence.

%%%%%%%%%%%%%%%%%%%%%%%%%%%%%%%%%%%%%%%%%%%%%%%%%%%%%%%%%%%%
\subsection{Reflection Principles}
\label{subsec:reflection-principles}
%%%%%%%%%%%%%%%%%%%%%%%%%%%%%%%%%%%%%%%%%%%%%%%%%%%%%%%%%%%%

A reflection principle states, informally, that whatever the theory proves
is true, at least for some specified class of formulas.

A schematic form is

\[
\boxed{
\operatorname{Pr}_T
(
\ulcorner\varphi\urcorner
)
\rightarrow
\varphi.
}
\]

Adding reflection principles generally increases proof-theoretic strength.

Full unrestricted self-reflection is constrained by incompleteness.

%%%%%%%%%%%%%%%%%%%%%%%%%%%%%%%%%%%%%%%%%%%%%%%%%%%%%%%%%%%%
\subsection{Derivability Conditions}
\label{subsec:derivability-conditions}
%%%%%%%%%%%%%%%%%%%%%%%%%%%%%%%%%%%%%%%%%%%%%%%%%%%%%%%%%%%%

Provability predicates for sufficiently strong formal theories satisfy
important internal properties.

Schematically, these include principles such as

\[
\boxed{
T\vdash\varphi
\Longrightarrow
T\vdash
\operatorname{Pr}_T(\ulcorner\varphi\urcorner),
}
\]

and internal closure of provability under modus ponens.

Such conditions are central in provability logic and incompleteness theory.

%%%%%%%%%%%%%%%%%%%%%%%%%%%%%%%%%%%%%%%%%%%%%%%%%%%%%%%%%%%%
\subsection{Intuitionistic Logic}
\label{subsec:proof-theory-intuitionistic-logic}
%%%%%%%%%%%%%%%%%%%%%%%%%%%%%%%%%%%%%%%%%%%%%%%%%%%%%%%%%%%%

Intuitionistic logic interprets proofs more constructively than classical
logic.

To prove

\[
\boxed{
A\lor B,
}
\]

one should provide a proof of \(A\) or a proof of \(B\).

To prove

\[
\boxed{
\exists x\,P(x),
}
\]

one should provide a witness \(a\) together with a proof of

\[
P(a).
\]

%%%%%%%%%%%%%%%%%%%%%%%%%%%%%%%%%%%%%%%%%%%%%%%%%%%%%%%%%%%%
\subsection{Excluded Middle}
\label{subsec:excluded-middle-proof-theory}
%%%%%%%%%%%%%%%%%%%%%%%%%%%%%%%%%%%%%%%%%%%%%%%%%%%%%%%%%%%%

Classical logic accepts

\[
\boxed{
P\lor\neg P
}
\]

for every proposition \(P\).

Intuitionistic logic does not accept this as a universal theorem without
additional information.

This does not mean that every instance is false.

It means a proof must constructively establish one side.

%%%%%%%%%%%%%%%%%%%%%%%%%%%%%%%%%%%%%%%%%%%%%%%%%%%%%%%%%%%%
\subsection{Double-Negation Elimination}
\label{subsec:double-negation-elimination-proof}
%%%%%%%%%%%%%%%%%%%%%%%%%%%%%%%%%%%%%%%%%%%%%%%%%%%%%%%%%%%%

Classical logic accepts

\[
\boxed{
\neg\neg P
\rightarrow
P.
}
\]

Intuitionistically,

\[
P
\rightarrow
\neg\neg P
\]

is valid, but the converse is not generally derivable.

This marks another important distinction between classical and constructive
proof systems.

%%%%%%%%%%%%%%%%%%%%%%%%%%%%%%%%%%%%%%%%%%%%%%%%%%%%%%%%%%%%
\subsection{Existence Proofs}
\label{subsec:constructive-existence-proofs}
%%%%%%%%%%%%%%%%%%%%%%%%%%%%%%%%%%%%%%%%%%%%%%%%%%%%%%%%%%%%

Classically, one may prove

\[
\exists x\,P(x)
\]

indirectly by showing that assuming no such \(x\) exists leads to
contradiction.

Constructively, a stronger expectation is often imposed:

\[
\boxed{
\text{produce an explicit witness }a
}
\]

and prove

\[
\boxed{
P(a).
}
\]

%%%%%%%%%%%%%%%%%%%%%%%%%%%%%%%%%%%%%%%%%%%%%%%%%%%%%%%%%%%%
\subsection{Curry--Howard Correspondence}
\label{subsec:curry-howard}
%%%%%%%%%%%%%%%%%%%%%%%%%%%%%%%%%%%%%%%%%%%%%%%%%%%%%%%%%%%%

The Curry--Howard correspondence identifies a deep analogy:

\[
\boxed{
\text{propositions}
\leftrightarrow
\text{types},
}
\]

and

\[
\boxed{
\text{proofs}
\leftrightarrow
\text{programs}.
}
\]

Under this perspective, constructing a proof is analogous to constructing
a program inhabiting a type.

%%%%%%%%%%%%%%%%%%%%%%%%%%%%%%%%%%%%%%%%%%%%%%%%%%%%%%%%%%%%
\subsection{Implication as Function Type}
\label{subsec:implication-function-type}
%%%%%%%%%%%%%%%%%%%%%%%%%%%%%%%%%%%%%%%%%%%%%%%%%%%%%%%%%%%%

A proof of

\[
A\rightarrow B
\]

takes a proof of \(A\) and transforms it into a proof of \(B\).

Thus implication corresponds naturally to a function type:

\[
\boxed{
A\rightarrow B.
}
\]

Proof application corresponds to function application.

%%%%%%%%%%%%%%%%%%%%%%%%%%%%%%%%%%%%%%%%%%%%%%%%%%%%%%%%%%%%
\subsection{Conjunction as Product Type}
\label{subsec:conjunction-product-type}
%%%%%%%%%%%%%%%%%%%%%%%%%%%%%%%%%%%%%%%%%%%%%%%%%%%%%%%%%%%%

A proof of

\[
A\land B
\]

contains:

\[
\boxed{
\text{a proof of }A
}
\]

and

\[
\boxed{
\text{a proof of }B.
}
\]

Thus conjunction corresponds to a product type.

The projections correspond to conjunction elimination.

%%%%%%%%%%%%%%%%%%%%%%%%%%%%%%%%%%%%%%%%%%%%%%%%%%%%%%%%%%%%
\subsection{Disjunction as Sum Type}
\label{subsec:disjunction-sum-type}
%%%%%%%%%%%%%%%%%%%%%%%%%%%%%%%%%%%%%%%%%%%%%%%%%%%%%%%%%%%%

A constructive proof of

\[
A\lor B
\]

contains either:

\[
\boxed{
\text{a proof of }A
}
\]

or

\[
\boxed{
\text{a proof of }B.
}
\]

Thus disjunction corresponds to a tagged sum type.

%%%%%%%%%%%%%%%%%%%%%%%%%%%%%%%%%%%%%%%%%%%%%%%%%%%%%%%%%%%%
\subsection{Existential Quantification as Dependent Pair}
\label{subsec:existential-dependent-pair}
%%%%%%%%%%%%%%%%%%%%%%%%%%%%%%%%%%%%%%%%%%%%%%%%%%%%%%%%%%%%

A constructive proof of

\[
\exists x\,P(x)
\]

contains:

\[
\boxed{
\text{a witness }a
}
\]

together with

\[
\boxed{
\text{a proof of }P(a).
}
\]

This resembles a dependent pair.

Type theory develops this correspondence in detail.

%%%%%%%%%%%%%%%%%%%%%%%%%%%%%%%%%%%%%%%%%%%%%%%%%%%%%%%%%%%%
\subsection{Normalization as Computation}
\label{subsec:normalization-as-computation}
%%%%%%%%%%%%%%%%%%%%%%%%%%%%%%%%%%%%%%%%%%%%%%%%%%%%%%%%%%%%

Under Curry--Howard, proof normalization corresponds to computation.

An introduction followed immediately by elimination resembles creating a
data structure and immediately destructuring it.

Removing the proof detour corresponds to reducing the associated program.

Thus

\[
\boxed{
\text{proof normalization}
\leftrightarrow
\text{program evaluation}.
}
\]

%%%%%%%%%%%%%%%%%%%%%%%%%%%%%%%%%%%%%%%%%%%%%%%%%%%%%%%%%%%%
\subsection{Proof Search}
\label{subsec:proof-theory-proof-search}
%%%%%%%%%%%%%%%%%%%%%%%%%%%%%%%%%%%%%%%%%%%%%%%%%%%%%%%%%%%%

Proof search attempts to construct a formal derivation automatically.

Given assumptions \(\Gamma\) and target \(\varphi\), the task is:

\[
\boxed{
\text{find a derivation of }
\Gamma\vdash\varphi.
}
\]

Search procedures use the structure of inference rules to determine which
subgoals to attempt.

%%%%%%%%%%%%%%%%%%%%%%%%%%%%%%%%%%%%%%%%%%%%%%%%%%%%%%%%%%%%
\subsection{Backward Proof Search}
\label{subsec:backward-proof-search}
%%%%%%%%%%%%%%%%%%%%%%%%%%%%%%%%%%%%%%%%%%%%%%%%%%%%%%%%%%%%

Backward search starts with the desired conclusion and asks which inference
rules could have produced it.

For example, to prove

\[
A\land B,
\]

conjunction introduction reduces the problem to two subgoals:

\[
\boxed{
A
}
\]

and

\[
\boxed{
B.
}
\]

Thus proof rules can function as algorithms for decomposing goals.

%%%%%%%%%%%%%%%%%%%%%%%%%%%%%%%%%%%%%%%%%%%%%%%%%%%%%%%%%%%%
\subsection{Forward Proof Search}
\label{subsec:forward-proof-search}
%%%%%%%%%%%%%%%%%%%%%%%%%%%%%%%%%%%%%%%%%%%%%%%%%%%%%%%%%%%%

Forward search begins with available assumptions and repeatedly applies
inference rules to generate consequences.

This can produce many irrelevant formulas.

Backward reasoning often focuses the search more directly on the goal.

Hybrid methods combine both directions.

%%%%%%%%%%%%%%%%%%%%%%%%%%%%%%%%%%%%%%%%%%%%%%%%%%%%%%%%%%%%
\subsection{Resolution}
\label{subsec:proof-theory-resolution}
%%%%%%%%%%%%%%%%%%%%%%%%%%%%%%%%%%%%%%%%%%%%%%%%%%%%%%%%%%%%

Resolution is a rule especially useful for automated reasoning.

From clauses

\[
\boxed{
A\lor C
}
\]

and

\[
\boxed{
\neg A\lor D,
}
\]

one may derive

\[
\boxed{
C\lor D.
}
\]

The complementary literals

\[
A
\]

and

\[
\neg A
\]

are eliminated.

%%%%%%%%%%%%%%%%%%%%%%%%%%%%%%%%%%%%%%%%%%%%%%%%%%%%%%%%%%%%
\subsection{Resolution Refutation}
\label{subsec:resolution-refutation}
%%%%%%%%%%%%%%%%%%%%%%%%%%%%%%%%%%%%%%%%%%%%%%%%%%%%%%%%%%%%

To prove formula \(\varphi\) by resolution:

\begin{enumerate}
    \item negate \(\varphi\);
    \item convert the result to clause form;
    \item repeatedly apply resolution;
    \item derive the empty clause.
\end{enumerate}

The empty clause represents contradiction.

Thus

\[
\boxed{
\neg\varphi
\text{ unsatisfiable}
\Longrightarrow
\varphi
\text{ valid}.
}
\]

%%%%%%%%%%%%%%%%%%%%%%%%%%%%%%%%%%%%%%%%%%%%%%%%%%%%%%%%%%%%
\subsection{Worked Example: Resolution}
\label{subsec:worked-resolution}
%%%%%%%%%%%%%%%%%%%%%%%%%%%%%%%%%%%%%%%%%%%%%%%%%%%%%%%%%%%%

\begin{workedexamplebox}{Resolving Two Clauses}

Suppose we have

\[
P\lor Q
\]

and

\[
\neg P\lor R.
\]

Resolve on \(P\).

The result is

\[
Q\lor R.
\]

\begin{answerbox}
\[
\boxed{
\frac{
P\lor Q
\qquad
\neg P\lor R
}{
Q\lor R
}.
}
\]
\end{answerbox}

\end{workedexamplebox}

%%%%%%%%%%%%%%%%%%%%%%%%%%%%%%%%%%%%%%%%%%%%%%%%%%%%%%%%%%%%
\subsection{Semantic Tableaux}
\label{subsec:semantic-tableaux}
%%%%%%%%%%%%%%%%%%%%%%%%%%%%%%%%%%%%%%%%%%%%%%%%%%%%%%%%%%%%

A tableau method systematically decomposes formulas into cases.

Branches represent candidate assignments or structures.

If every branch closes in contradiction, the original set of formulas is
unsatisfiable.

Tableaux therefore combine proof search with semantic intuition.

%%%%%%%%%%%%%%%%%%%%%%%%%%%%%%%%%%%%%%%%%%%%%%%%%%%%%%%%%%%%
\subsection{Closed Tableau}
\label{subsec:closed-tableau}
%%%%%%%%%%%%%%%%%%%%%%%%%%%%%%%%%%%%%%%%%%%%%%%%%%%%%%%%%%%%

A branch closes when it contains both

\[
\boxed{
A
}
\]

and

\[
\boxed{
\neg A.
}
\]

If all branches close, no consistent interpretation survives.

Therefore the starting assumptions are inconsistent.

%%%%%%%%%%%%%%%%%%%%%%%%%%%%%%%%%%%%%%%%%%%%%%%%%%%%%%%%%%%%
\subsection{Proof Complexity}
\label{subsec:proof-complexity}
%%%%%%%%%%%%%%%%%%%%%%%%%%%%%%%%%%%%%%%%%%%%%%%%%%%%%%%%%%%%

Proof complexity studies the resources required to prove statements in
formal systems.

Possible measures include:

\[
\boxed{
\text{proof length},
}
\]

\[
\boxed{
\text{proof depth},
}
\]

and

\[
\boxed{
\text{formula size}.
}
\]

Two proof systems may prove the same theorems but require radically
different proof lengths.

%%%%%%%%%%%%%%%%%%%%%%%%%%%%%%%%%%%%%%%%%%%%%%%%%%%%%%%%%%%%
\subsection{Polynomial Simulation of Proof Systems}
\label{subsec:proof-system-simulation}
%%%%%%%%%%%%%%%%%%%%%%%%%%%%%%%%%%%%%%%%%%%%%%%%%%%%%%%%%%%%

One proof system simulates another efficiently if proofs in the second can
be transformed into proofs in the first with only controlled growth.

If proof size increases polynomially, one speaks of polynomial simulation.

This allows formal comparison of proof-system efficiency.

%%%%%%%%%%%%%%%%%%%%%%%%%%%%%%%%%%%%%%%%%%%%%%%%%%%%%%%%%%%%
\subsection{Proof Length Versus Truth}
\label{subsec:proof-length-versus-truth}
%%%%%%%%%%%%%%%%%%%%%%%%%%%%%%%%%%%%%%%%%%%%%%%%%%%%%%%%%%%%

The existence of a proof does not imply the existence of a short proof.

A theorem may be derivable while every proof in a particular proof system
is extremely long.

Thus proof complexity studies an additional dimension beyond mere
provability.

%%%%%%%%%%%%%%%%%%%%%%%%%%%%%%%%%%%%%%%%%%%%%%%%%%%%%%%%%%%%
\subsection{Proof Checking}
\label{subsec:formal-proof-checking}
%%%%%%%%%%%%%%%%%%%%%%%%%%%%%%%%%%%%%%%%%%%%%%%%%%%%%%%%%%%%

Given a purported formal proof, verification is usually local.

Each line can be checked against:

\[
\boxed{
\text{the axioms},
}
\]

and

\[
\boxed{
\text{the inference rules}.
}
\]

This is much easier than discovering the proof in the first place.

%%%%%%%%%%%%%%%%%%%%%%%%%%%%%%%%%%%%%%%%%%%%%%%%%%%%%%%%%%%%
\subsection{Proof Certificates}
\label{subsec:proof-certificates}
%%%%%%%%%%%%%%%%%%%%%%%%%%%%%%%%%%%%%%%%%%%%%%%%%%%%%%%%%%%%

A proof certificate is a formal object that allows an independent checker
to verify a theorem.

The checker can be small and trusted even if the proof was generated by a
large or complex automated system.

This separation is important in formal verification.

%%%%%%%%%%%%%%%%%%%%%%%%%%%%%%%%%%%%%%%%%%%%%%%%%%%%%%%%%%%%
\subsection{Automated Theorem Proving}
\label{subsec:automated-theorem-proving-proof-theory}
%%%%%%%%%%%%%%%%%%%%%%%%%%%%%%%%%%%%%%%%%%%%%%%%%%%%%%%%%%%%

Automated theorem provers search formal proof spaces using methods such as:

\[
\boxed{
\text{resolution},
}
\]

\[
\boxed{
\text{rewriting},
}
\]

\[
\boxed{
\text{tableaux},
}
\]

and

\[
\boxed{
\text{sequent search}.
}
\]

Proof theory provides the formal systems and correctness principles behind
these tools.

%%%%%%%%%%%%%%%%%%%%%%%%%%%%%%%%%%%%%%%%%%%%%%%%%%%%%%%%%%%%
\subsection{Interactive Theorem Proving}
\label{subsec:interactive-theorem-proving}
%%%%%%%%%%%%%%%%%%%%%%%%%%%%%%%%%%%%%%%%%%%%%%%%%%%%%%%%%%%%

In interactive theorem proving, a human guides the proof while software
checks every formal step.

The workflow is roughly

\[
\boxed{
\text{human strategy}
+
\text{machine verification}.
}
\]

This is useful when automated proof search alone is too difficult.

%%%%%%%%%%%%%%%%%%%%%%%%%%%%%%%%%%%%%%%%%%%%%%%%%%%%%%%%%%%%
\subsection{Proof Assistants and Trusted Kernels}
\label{subsec:proof-assistants-kernels}
%%%%%%%%%%%%%%%%%%%%%%%%%%%%%%%%%%%%%%%%%%%%%%%%%%%%%%%%%%%%

Many proof assistants use a small trusted kernel.

Higher-level automation proposes proof terms or derivations.

The kernel checks whether those objects satisfy the formal inference rules.

Thus confidence is concentrated in a relatively small verification
component.

%%%%%%%%%%%%%%%%%%%%%%%%%%%%%%%%%%%%%%%%%%%%%%%%%%%%%%%%%%%%
\subsection{Proof Theory and Computability}
\label{subsec:proof-theory-computability}
%%%%%%%%%%%%%%%%%%%%%%%%%%%%%%%%%%%%%%%%%%%%%%%%%%%%%%%%%%%%

Proof systems and computation are deeply connected.

Important questions include:

\[
\boxed{
\text{Can proofs be effectively enumerated?}
}
\]

\[
\boxed{
\text{Is theoremhood decidable?}
}
\]

and

\[
\boxed{
\text{Can proof search always terminate?}
}
\]

These lead directly to computability theory.

%%%%%%%%%%%%%%%%%%%%%%%%%%%%%%%%%%%%%%%%%%%%%%%%%%%%%%%%%%%%
\subsection{Proof Theory and Type Theory}
\label{subsec:proof-theory-type-theory}
%%%%%%%%%%%%%%%%%%%%%%%%%%%%%%%%%%%%%%%%%%%%%%%%%%%%%%%%%%%%

Through Curry--Howard,

\[
\boxed{
\text{proof theory}
}
\]

and

\[
\boxed{
\text{type theory}
}
\]

become closely connected.

Logical inference corresponds to construction of typed terms.

Normalization corresponds to program reduction.

This connection is developed further in Section~12.7.

%%%%%%%%%%%%%%%%%%%%%%%%%%%%%%%%%%%%%%%%%%%%%%%%%%%%%%%%%%%%
\subsection{Proof Theory and Constructive Mathematics}
\label{subsec:proof-theory-constructive-mathematics}
%%%%%%%%%%%%%%%%%%%%%%%%%%%%%%%%%%%%%%%%%%%%%%%%%%%%%%%%%%%%

Constructive mathematics pays close attention to what information a proof
provides.

A proof of existence may be required to construct a witness.

A proof of disjunction may be required to determine which disjunct holds.

Proof theory provides formal systems for analyzing these constructive
requirements.

%%%%%%%%%%%%%%%%%%%%%%%%%%%%%%%%%%%%%%%%%%%%%%%%%%%%%%%%%%%%
\subsection{Proof-Theoretic Workflow}
\label{subsec:proof-theory-workflow}
%%%%%%%%%%%%%%%%%%%%%%%%%%%%%%%%%%%%%%%%%%%%%%%%%%%%%%%%%%%%

When analyzing a formal proof system:

\begin{enumerate}
    \item specify the formal language;
    \item identify logical and nonlogical axioms;
    \item state the inference rules;
    \item define what counts as a derivation;
    \item distinguish assumptions from discharged assumptions;
    \item determine whether the system is Hilbert-style, natural
          deduction, sequent calculus, or another formalism;
    \item prove soundness of the inference rules;
    \item investigate completeness where appropriate;
    \item search for normal forms of proofs;
    \item identify introduction-elimination detours;
    \item investigate cut elimination;
    \item use the subformula property when available;
    \item analyze consistency;
    \item compare the strength of theories;
    \item determine whether extensions are conservative;
    \item encode syntax arithmetically when studying self-reference;
    \item distinguish logical completeness from theory completeness;
    \item state the hypotheses of incompleteness results carefully;
    \item use ordinal measures when analyzing transfinite proof reduction;
    \item study proof search and proof length separately from mere
          derivability;
    \item verify formal proofs with independent checking when possible.
\end{enumerate}

%%%%%%%%%%%%%%%%%%%%%%%%%%%%%%%%%%%%%%%%%%%%%%%%%%%%%%%%%%%%
\subsection{Common Errors}
\label{subsec:proof-theory-common-errors}
%%%%%%%%%%%%%%%%%%%%%%%%%%%%%%%%%%%%%%%%%%%%%%%%%%%%%%%%%%%%

\subsubsection{Confusing Proof Theory With Model Theory}

Proof theory primarily studies

\[
\boxed{
\text{formal derivations}.
}
\]

Model theory primarily studies

\[
\boxed{
\text{structures satisfying theories}.
}
\]

The two are linked by soundness and completeness.

%%%%%%%%%%%%%%%%%%%%%%%%%%%%%%%%%%%%%%%%%%%%%%%%%%%%%%%%%%%%
\subsubsection{Confusing \(\vdash\) and \(\models\)}

The notation

\[
\Gamma\vdash\varphi
\]

is syntactic.

The notation

\[
\Gamma\models\varphi
\]

is semantic.

Proof theory concentrates primarily on the former.

%%%%%%%%%%%%%%%%%%%%%%%%%%%%%%%%%%%%%%%%%%%%%%%%%%%%%%%%%%%%
\subsubsection{Assuming Every Proof System Uses the Same Rules}

Hilbert systems, natural deduction, sequent calculus, tableaux, and
resolution have different proof structures.

They may nevertheless characterize the same logical consequence relation.

%%%%%%%%%%%%%%%%%%%%%%%%%%%%%%%%%%%%%%%%%%%%%%%%%%%%%%%%%%%%
\subsubsection{Assuming Cut Elimination Means Lemmas Are Useless}

Cut elimination means intermediate lemmas are not necessary for
derivability in the formal system.

Human mathematicians still use lemmas because they can drastically shorten
and clarify proofs.

%%%%%%%%%%%%%%%%%%%%%%%%%%%%%%%%%%%%%%%%%%%%%%%%%%%%%%%%%%%%
\subsubsection{Assuming Cut-Free Proofs Are Always Shorter}

Eliminating cuts can cause a proof to become much larger.

Structural simplicity and proof length are different issues.

%%%%%%%%%%%%%%%%%%%%%%%%%%%%%%%%%%%%%%%%%%%%%%%%%%%%%%%%%%%%
\subsubsection{Confusing Normalization With Semantic Simplification}

Normalization transforms proof structure.

It does not merely simplify the written formula being proved.

%%%%%%%%%%%%%%%%%%%%%%%%%%%%%%%%%%%%%%%%%%%%%%%%%%%%%%%%%%%%
\subsubsection{Assuming Consistency Means Completeness}

A consistent theory avoids contradiction.

A complete theory decides every sentence.

A theory may be consistent but incomplete.

%%%%%%%%%%%%%%%%%%%%%%%%%%%%%%%%%%%%%%%%%%%%%%%%%%%%%%%%%%%%
\subsubsection{Confusing Relative Consistency With Proof of Consistency}

A result

\[
\operatorname{Con}(T)
\rightarrow
\operatorname{Con}(T')
\]

does not establish \(\operatorname{Con}(T)\).

It only compares the theories.

%%%%%%%%%%%%%%%%%%%%%%%%%%%%%%%%%%%%%%%%%%%%%%%%%%%%%%%%%%%%
\subsubsection{Confusing Gödel Completeness and Incompleteness}

First-order logical completeness concerns the relation between

\[
\models
\]

and

\[
\vdash.
\]

Gödel incompleteness concerns whether a particular sufficiently strong
theory decides every sentence.

These are different questions.

%%%%%%%%%%%%%%%%%%%%%%%%%%%%%%%%%%%%%%%%%%%%%%%%%%%%%%%%%%%%
\subsubsection{Assuming Every Consistent Theory Is Incomplete}

Incompleteness requires specific hypotheses, including sufficient
arithmetical expressive strength and effective axiomatizability.

Many weaker theories are complete.

%%%%%%%%%%%%%%%%%%%%%%%%%%%%%%%%%%%%%%%%%%%%%%%%%%%%%%%%%%%%
\subsubsection{Saying a Gödel Sentence Is Simply ``This Sentence Is False''}

Gödel's construction is not the liar paradox.

The sentence is built through formal arithmetization and provability.

Its behavior depends on the formal theory and the representation of
provability.

%%%%%%%%%%%%%%%%%%%%%%%%%%%%%%%%%%%%%%%%%%%%%%%%%%%%%%%%%%%%
\subsubsection{Assuming Incompleteness Makes Formal Proof Meaningless}

Incompleteness limits what one formal theory can prove.

It does not undermine ordinary valid deductions inside the theory.

%%%%%%%%%%%%%%%%%%%%%%%%%%%%%%%%%%%%%%%%%%%%%%%%%%%%%%%%%%%%
\subsubsection{Assuming a Theory Can Prove Its Own Consistency Because It Is Consistent}

Gödel's Second Incompleteness Theorem blocks this for sufficiently strong
effective consistent theories under the standard formalization of
consistency.

%%%%%%%%%%%%%%%%%%%%%%%%%%%%%%%%%%%%%%%%%%%%%%%%%%%%%%%%%%%%
\subsubsection{Assuming Intuitionistic Logic Declares Excluded Middle False}

Intuitionistic logic does not accept excluded middle as a universal theorem.

Specific instances may still be provable.

%%%%%%%%%%%%%%%%%%%%%%%%%%%%%%%%%%%%%%%%%%%%%%%%%%%%%%%%%%%%
\subsubsection{Confusing Constructive Existence With Numerical Approximation}

Constructive existence means a proof provides sufficient information to
construct a witness in the relevant formal sense.

It is not merely an approximate numerical answer.

%%%%%%%%%%%%%%%%%%%%%%%%%%%%%%%%%%%%%%%%%%%%%%%%%%%%%%%%%%%%
\subsubsection{Confusing Curry--Howard With an Identity of All Logic and All Programming}

Curry--Howard gives a structural correspondence between suitable proof
systems and typed calculi.

Its exact form depends on the logical and type-theoretic systems involved.

%%%%%%%%%%%%%%%%%%%%%%%%%%%%%%%%%%%%%%%%%%%%%%%%%%%%%%%%%%%%
\subsubsection{Assuming Proof Search and Proof Checking Have the Same Difficulty}

Checking a supplied proof is usually local and straightforward.

Finding a proof can require enormous search.

%%%%%%%%%%%%%%%%%%%%%%%%%%%%%%%%%%%%%%%%%%%%%%%%%%%%%%%%%%%%
\subsubsection{Assuming Short Statements Have Short Proofs}

Formula length and proof length need not correlate.

Some short statements can require long proofs in a particular proof system.

%%%%%%%%%%%%%%%%%%%%%%%%%%%%%%%%%%%%%%%%%%%%%%%%%%%%%%%%%%%%
\subsection{Applications}
\label{subsec:proof-theory-applications}
%%%%%%%%%%%%%%%%%%%%%%%%%%%%%%%%%%%%%%%%%%%%%%%%%%%%%%%%%%%%

\begin{applicationbox}

\textbf{Foundations of Mathematics.}
Proof theory studies consistency, formal derivability, incompleteness, and
the relative strength of axiomatic systems.

\medskip

\textbf{Formal Verification.}
Proof rules provide the basis for mechanically checked proofs of software
and hardware correctness.

\medskip

\textbf{Automated Theorem Proving.}
Resolution, tableaux, sequent calculus, and proof search are used by
automated reasoning systems.

\medskip

\textbf{Type Theory.}
The Curry--Howard correspondence identifies proofs with typed terms and
normalization with computation.

\medskip

\textbf{Constructive Mathematics.}
Proof theory formalizes mathematical systems in which existence and
disjunction carry computational information.

\medskip

\textbf{Computability.}
Arithmetization of syntax and provability connects formal proof with
algorithmic decidability and undecidability.

\medskip

\textbf{Proof Complexity.}
Formal proof systems can be compared by the size and structure of proofs
they require.

\end{applicationbox}

%%%%%%%%%%%%%%%%%%%%%%%%%%%%%%%%%%%%%%%%%%%%%%%%%%%%%%%%%%%%
\subsection{Exercises}
\label{subsec:proof-theory-exercises}
%%%%%%%%%%%%%%%%%%%%%%%%%%%%%%%%%%%%%%%%%%%%%%%%%%%%%%%%%%%%

\begin{exercise}[1]
Define proof theory.
\end{exercise}

\begin{exercise}[1]
Define a formal derivation.
\end{exercise}

\begin{exercise}[1]
Define an inference rule.
\end{exercise}

\begin{exercise}[1]
Define an axiom schema.
\end{exercise}

\begin{exercise}[1]
Define a sequent.
\end{exercise}

\begin{exercise}[1]
Define cut elimination.
\end{exercise}

\begin{exercise}[1]
Define normalization.
\end{exercise}

\begin{exercise}[1]
Define proof-theoretic consistency.
\end{exercise}

\begin{exercise}[1]
Define a conservative extension.
\end{exercise}

\begin{exercise}[1]
Define proof-theoretic strength conceptually.
\end{exercise}

\begin{exercise}[2]
Construct a natural-deduction derivation of

\[
P\land Q\rightarrow P.
\]
\end{exercise}

\begin{exercise}[2]
Construct a derivation of

\[
P\rightarrow(Q\rightarrow P).
\]
\end{exercise}

\begin{exercise}[2]
Use conjunction introduction to derive

\[
P\land Q
\]

from \(P\) and \(Q\).
\end{exercise}

\begin{exercise}[2]
Use disjunction introduction to derive

\[
P\lor Q
\]

from \(P\).
\end{exercise}

\begin{exercise}[3]
Construct a proof by cases using disjunction elimination.
\end{exercise}

\begin{exercise}[3]
Explain assumption discharge in implication introduction.
\end{exercise}

\begin{exercise}[3]
Explain proof by contradiction formally.
\end{exercise}

\begin{exercise}[2]
Interpret the sequent

\[
P,Q\Rightarrow R.
\]
\end{exercise}

\begin{exercise}[2]
Give an example of left weakening.
\end{exercise}

\begin{exercise}[2]
Give an example of contraction.
\end{exercise}

\begin{exercise}[2]
Give an example of exchange.
\end{exercise}

\begin{exercise}[3]
Explain the logical meaning of the cut rule.
\end{exercise}

\begin{exercise}[3]
Explain why cut resembles use of an intermediate lemma.
\end{exercise}

\begin{exercise}[3]
State Gentzen's cut-elimination theorem conceptually.
\end{exercise}

\begin{exercise}[3]
Explain the subformula property.
\end{exercise}

\begin{exercise}[3]
Explain why the subformula property aids proof search.
\end{exercise}

\begin{exercise}[3]
Give an example of an introduction-elimination detour.
\end{exercise}

\begin{exercise}[3]
Normalize that proof.
\end{exercise}

\begin{exercise}[3]
Explain how normalization can establish consistency.
\end{exercise}

\begin{exercise}[2]
Distinguish syntactic consistency and semantic satisfiability.
\end{exercise}

\begin{exercise}[3]
Explain how first-order completeness connects the two.
\end{exercise}

\begin{exercise}[3]
Give an example of a conservative extension conceptually.
\end{exercise}

\begin{exercise}[3]
Explain relative consistency.
\end{exercise}

\begin{exercise}[3]
Describe the induction schema of first-order arithmetic.
\end{exercise}

\begin{exercise}[3]
Explain how restricting induction changes theory strength.
\end{exercise}

\begin{exercise}[2]
Explain Gödel coding conceptually.
\end{exercise}

\begin{exercise}[3]
Explain arithmetization of syntax.
\end{exercise}

\begin{exercise}[3]
Describe the meaning of

\[
\operatorname{Proof}_T(p,n).
\]
\end{exercise}

\begin{exercise}[3]
Describe the meaning of

\[
\operatorname{Pr}_T(n).
\]
\end{exercise}

\begin{exercise}[3]
Explain the purpose of the diagonal lemma.
\end{exercise}

\begin{exercise}[3]
State Gödel's First Incompleteness Theorem conceptually.
\end{exercise}

\begin{exercise}[3]
State Gödel's Second Incompleteness Theorem conceptually.
\end{exercise}

\begin{exercise}[3]
Explain why first-order completeness does not contradict incompleteness.
\end{exercise}

\begin{exercise}[3]
Explain why incompleteness does not apply to every formal theory.
\end{exercise}

\begin{exercise}[3]
Explain the role of effective axiomatizability.
\end{exercise}

\begin{exercise}[3]
Explain the formal consistency statement

\[
\operatorname{Con}(T).
\]
\end{exercise}

\begin{exercise}[3]
Describe Hilbert's program.
\end{exercise}

\begin{exercise}[3]
Explain how Gödel's theorems affect Hilbert's program.
\end{exercise}

\begin{exercise}[3]
Explain Gentzen's consistency strategy conceptually.
\end{exercise}

\begin{exercise}[3]
Define a proof-theoretic ordinal conceptually.
\end{exercise}

\begin{exercise}[3]
Explain the role of \(\varepsilon_0\).
\end{exercise}

\begin{exercise}[3]
Explain ordinal descent in a proof reduction argument.
\end{exercise}

\begin{exercise}[3]
Explain a reflection principle conceptually.
\end{exercise}

\begin{exercise}[2]
State the law of excluded middle.
\end{exercise}

\begin{exercise}[2]
State double-negation elimination.
\end{exercise}

\begin{exercise}[3]
Distinguish classical and intuitionistic logic.
\end{exercise}

\begin{exercise}[3]
Explain what constructive proof of

\[
A\lor B
\]

requires.
\end{exercise}

\begin{exercise}[3]
Explain what constructive proof of

\[
\exists x\,P(x)
\]

requires.
\end{exercise}

\begin{exercise}[3]
Describe the Curry--Howard correspondence.
\end{exercise}

\begin{exercise}[3]
Explain why implication corresponds to a function type.
\end{exercise}

\begin{exercise}[3]
Explain why conjunction corresponds to a product type.
\end{exercise}

\begin{exercise}[3]
Explain why disjunction corresponds to a sum type.
\end{exercise}

\begin{exercise}[3]
Explain normalization as computation.
\end{exercise}

\begin{exercise}[2]
Perform one resolution step on

\[
P\lor Q
\]

and

\[
\neg P\lor R.
\]
\end{exercise}

\begin{exercise}[3]
Explain resolution refutation.
\end{exercise}

\begin{exercise}[3]
Explain how a semantic tableau branch closes.
\end{exercise}

\begin{exercise}[3]
Compare backward and forward proof search.
\end{exercise}

\begin{exercise}[3]
Explain proof complexity.
\end{exercise}

\begin{exercise}[3]
Distinguish proof length from statement length.
\end{exercise}

\begin{exercise}[3]
Explain polynomial simulation between proof systems.
\end{exercise}

\begin{exercise}[3]
Explain why proof checking is usually easier than proof search.
\end{exercise}

\begin{exercise}[3]
Explain the purpose of a proof certificate.
\end{exercise}

\begin{exercise}[3]
Explain the role of a trusted kernel in a proof assistant.
\end{exercise}

\begin{exercise}[4]
Study a Hilbert-style axiomatization of propositional logic.
\end{exercise}

\begin{exercise}[4]
Prove soundness of natural deduction for propositional logic.
\end{exercise}

\begin{exercise}[4]
Study completeness of natural deduction.
\end{exercise}

\begin{exercise}[4]
Develop a classical sequent calculus.
\end{exercise}

\begin{exercise}[4]
Develop an intuitionistic sequent calculus.
\end{exercise}

\begin{exercise}[4]
Study the proof of cut elimination.
\end{exercise}

\begin{exercise}[4]
Study the relationship between cut elimination and consistency.
\end{exercise}

\begin{exercise}[4]
Study normalization for natural deduction.
\end{exercise}

\begin{exercise}[4]
Study strong normalization.
\end{exercise}

\begin{exercise}[4]
Study the Church--Rosser property conceptually.
\end{exercise}

\begin{exercise}[4]
Study Gödel numbering in detail.
\end{exercise}

\begin{exercise}[4]
Study representability of recursive relations in arithmetic.
\end{exercise}

\begin{exercise}[4]
Study the diagonal lemma.
\end{exercise}

\begin{exercise}[4]
Study the proof of Gödel's First Incompleteness Theorem.
\end{exercise}

\begin{exercise}[4]
Study Rosser's strengthening of the First Incompleteness Theorem.
\end{exercise}

\begin{exercise}[4]
Study the derivability conditions for provability.
\end{exercise}

\begin{exercise}[4]
Study Gödel's Second Incompleteness Theorem.
\end{exercise}

\begin{exercise}[4]
Study Löb's theorem.
\end{exercise}

\begin{exercise}[4]
Study provability logic.
\end{exercise}

\begin{exercise}[4]
Study Gentzen's consistency proof for arithmetic.
\end{exercise}

\begin{exercise}[4]
Study ordinal notation below \(\varepsilon_0\).
\end{exercise}

\begin{exercise}[4]
Study proof-theoretic ordinal analysis.
\end{exercise}

\begin{exercise}[4]
Study intuitionistic natural deduction.
\end{exercise}

\begin{exercise}[4]
Study the Curry--Howard correspondence formally.
\end{exercise}

\begin{exercise}[4]
Study resolution completeness.
\end{exercise}

\begin{exercise}[4]
Study proof complexity of propositional systems.
\end{exercise}

\begin{exercise}[4]
Compare Frege systems, resolution, and sequent systems.
\end{exercise}

%%%%%%%%%%%%%%%%%%%%%%%%%%%%%%%%%%%%%%%%%%%%%%%%%%%%%%%%%%%%
\subsection{Conceptual Summary}
\label{subsec:proof-theory-summary}
%%%%%%%%%%%%%%%%%%%%%%%%%%%%%%%%%%%%%%%%%%%%%%%%%%%%%%%%%%%%

Proof theory studies

\[
\boxed{
\text{formal proofs as mathematical objects}.
}
\]

A formal derivation is built from:

\[
\boxed{
\text{axioms}
+
\text{assumptions}
+
\text{inference rules}.
}
\]

Major proof systems include:

\[
\boxed{
\text{Hilbert systems},
}
\]

\[
\boxed{
\text{natural deduction},
}
\]

and

\[
\boxed{
\text{sequent calculus}.
}
\]

Sequent calculus introduces the cut rule:

\[
\boxed{
\frac{
\Gamma\Rightarrow\Delta,A
\qquad
A,\Pi\Rightarrow\Lambda
}{
\Gamma,\Pi\Rightarrow\Delta,\Lambda
}.
}
\]

Cut elimination shows that many such intermediate formulas can be removed.

Natural-deduction normalization similarly removes unnecessary proof
detours.

Proof theory also studies

\[
\boxed{
\text{consistency},
}
\]

\[
\boxed{
\text{proof-theoretic strength},
}
\]

\[
\boxed{
\text{formalized provability},
}
\]

and

\[
\boxed{
\text{incompleteness}.
}
\]

Gödel coding makes it possible for arithmetic to represent statements
about its own formulas and proofs.

This leads to fundamental limitations on sufficiently strong effective
formal theories.

At the constructive end, Curry--Howard connects

\[
\boxed{
\text{propositions with types}
}
\]

and

\[
\boxed{
\text{proofs with programs}.
}
\]

Proof theory therefore forms a bridge among

\[
\boxed{
\text{logic}
+
\text{computation}
+
\text{foundations}
+
\text{formal verification}.
}
\]

%%%%%%%%%%%%%%%%%%%%%%%%%%%%%%%%%%%%%%%%%%%%%%%%%%%%%%%%%%%%
\subsection{Section Connections}
\label{subsec:proof-theory-connections}
%%%%%%%%%%%%%%%%%%%%%%%%%%%%%%%%%%%%%%%%%%%%%%%%%%%%%%%%%%%%

\chapterconnection{
Proof Theory develops the syntactic counterpart to the semantic viewpoint
of Section~\ref{sec:model-theory}. Formal Logic introduced derivability,
soundness, and completeness, while Proof Theory studies the internal
structure of derivations through natural deduction, sequent calculus,
normalization, cut elimination, consistency, and proof-theoretic strength.
Gödel coding and formal provability connect proofs with arithmetic and lead
to the incompleteness theorems. The next section, Computability Theory,
develops the algorithmic side of these ideas by formalizing effective
procedures, computable functions, Turing machines, decidability, and
undecidability. Later sections on Recursion Theory, Type Theory, and
Constructive Mathematics deepen the connection between proof and
computation.
}

%%%%%%%%%%%%%%%%%%%%%%%%%%%%%%%%%%%%%%%%%%%%%%%%%%%%%%%%%%%%
% SECTION SOURCE TRACKING
%%%%%%%%%%%%%%%%%%%%%%%%%%%%%%%%%%%%%%%%%%%%%%%%%%%%%%%%%%%%

%------------------------------------------------------------
% 12.4 Proof Theory
%------------------------------------------------------------
%
% Source basis:
%   - Standard proof theory.
%   - Standard natural deduction and sequent calculus.
%   - Standard normalization, cut-elimination, incompleteness,
%     ordinal-analysis, and constructive-proof concepts.
%
% Used for:
%   - proofs as mathematical objects;
%   - object theory and metatheory;
%   - axioms and inference rules;
%   - derivations and proof trees;
%   - Hilbert-style systems;
%   - axiom schemas;
%   - natural deduction;
%   - conjunction, implication, disjunction, and negation rules;
%   - classical reasoning;
%   - sequent calculus;
%   - structural rules;
%   - weakening, contraction, and exchange;
%   - left and right rules;
%   - cut;
%   - cut elimination;
%   - subformula property;
%   - proof search;
%   - normalization;
%   - introduction-elimination detours;
%   - consistency;
%   - syntactic versus semantic consistency;
%   - consistency from cut elimination;
%   - conservative extensions;
%   - relative consistency;
%   - proof-theoretic strength;
%   - formal arithmetic;
%   - induction schemas;
%   - Gödel coding;
%   - arithmetization of syntax;
%   - provability predicates;
%   - self-reference;
%   - diagonal lemma;
%   - First Incompleteness Theorem;
%   - logical completeness versus incompleteness;
%   - Second Incompleteness Theorem;
%   - formal consistency statements;
%   - Hilbert's program;
%   - finitistic reasoning;
%   - Gentzen's consistency program;
%   - proof-theoretic ordinals;
%   - epsilon-zero;
%   - ordinal descent;
%   - reflection principles;
%   - derivability conditions;
%   - intuitionistic logic;
%   - excluded middle;
%   - double-negation elimination;
%   - constructive existence;
%   - Curry--Howard correspondence;
%   - implication, product, sum, and dependent-pair interpretations;
%   - normalization as computation;
%   - proof search;
%   - resolution;
%   - semantic tableaux;
%   - proof complexity;
%   - proof checking;
%   - proof certificates;
%   - automated and interactive theorem proving;
%   - applications and exercises.
%
% Notes:
%   - Computability Theory is developed next in Section 12.5.
%   - Formal Logic appears in Section 12.1.
%   - Model Theory appears in Section 12.3.
%   - Recursion Theory follows in Section 12.6.
%   - Type Theory follows in Section 12.7.
%
%%%%%%%%%%%%%%%%%%%%%%%%%%%%%%%%%%%%%%%%%%%%%%%%%%%%%%%%%%%%